\section{Deletion and projective normalization}\label{alt:section}

The preceding estimates allow us to remove every repeated projected point
and every small triangle while retaining many columns. We first construct
a set for each sufficiently large prime $r$, compare its area bound with
its cardinality, and then pass to every sufficiently large cardinality.
All constants below may depend on the fixed integer $k$; none depends on
the growing prime $r$.

This final step is an elementary alteration argument. The earlier
probabilistic lower bound of Koml\'os, Pintz and
Szemer\'edi~\cite{KPS1982} uses a stronger hypergraph independence lemma;
no such lemma is needed below.

\begin{proof}[Proof of Theorem~\ref{intro:main}]
Keep the parameters constructed above, with $d=41$, and set
\begin{equation}\label{alt:sample-size}
 a_r=r\sqrt{N^3/\tau},\qquad n_r=\lfloor a_r\rfloor.
\end{equation}
Here $k\ge2$ and $B$ is an odd prime, so $\tau\ge1$.  Take $2n_r$ samples
with the shared random choices and independent conditional draws specified
in Section~\ref{sam:section}.  Let $Z$ count the unordered pairs of sample
indices with equal projected points, plus the unordered triples whose
projected points are pairwise distinct and whose determinant has absolute
value at most $\tau$.  Linearity of expectation, Lemma~\ref{sam:proportional},
and Corollary~\ref{sam:bad-triple} give
\begin{equation}\label{alt:expected-events}
 \frac{\mathbb EZ}{n_r}
 \ll n_r(hq)^6N^{-3}
    +n_r^2(\log(2N))^2\frac{h}{N^3r^d}.
\end{equation}
No independence between the counted events is needed.

Both terms tend to zero.  Indeed, $N=(hq)^{10}$ and $\tau\ge1$ give
\[
 n_r(hq)^6N^{-3}\le r\tau^{-1/2}(hq)^{-9}=o(1).
\]
For the other term, the definition of $n_r$ gives
\[
 n_r^2(\log(2N))^2\frac{h}{N^3r^d}
 \le (\log(2N))^2\frac{h}{\tau}r^{2-d}.
\]
Since $B^{k-1}\ge3$, the definition of $\tau$ implies
$\tau\ge B^{k-1}/3$, and hence $h/\tau\le3B=O_k(r^{30})$.
The bounds $B=O_k(r^{30})$, $q\le2h^{100}$, and $N=(hq)^{10}$
also give $\log(2N)=O_k(\log r)$. Thus the second term in
\eqref{alt:expected-events} is
$O_k((\log r)^2r^{-9})=o(1)$.
Thus $\mathbb EZ=o(n_r)$.

For sufficiently large $r$, choose an outcome with $Z<n_r$.  For each
violating pair or triple in that outcome, mark one of its indices for
deletion.  This marks at most $Z$ indices and meets every original
violation.  After deleting all marked indices, at least $n_r$ remain;
retain any $n_r$ of them.  Their projections form a set $P_r$ of $n_r$ distinct
points in $[0,1)^2$, and every triple of retained columns satisfies
$|\det A|>\tau$.  In particular, the construction excludes every
zero-area triple as well.

For a retained triple, divide the $j$th column of $A$ by its positive
third coordinate $A_{3j}$.  The resulting columns have the form
$(p_{j1},p_{j2},1)^{\mathsf T}$, so the determinant formula for triangle
area gives the exact identity
\begin{equation}\label{alt:area-identity}
 \operatorname{Area}(p_1p_2p_3)
 =\frac{|\det A|}{2A_{31}A_{32}A_{33}}.
\end{equation}
Since $N\le A_{3j}<2N$ for each $j$, every retained triangle has area at
least $\tau/(16N^3)$.  Consequently
\begin{equation}\label{alt:area-bound}
 \Delta(P_r)\ge\frac{\tau}{16N^3}.
\end{equation}
Also $\tau<h$ and $N\ge h^{10}$ give
$a_r>rh^{29/2}\to\infty$.  Hence $n_r\to\infty$ and eventually
$n_r\ge a_r/2$.  Combining this with \eqref{alt:area-bound} yields
\begin{equation}\label{alt:scaled-area}
 n_r^2\Delta(P_r)\ge\frac{r^2}{64}.
\end{equation}

We now express both the cardinality and the area bound in terms of $r$.
The parameter choices give
\[
 \begin{gathered}
  100k^2r^{30}<B\le200k^2r^{30},\qquad h=B^k,\\
  h^{100}<q\le2h^{100},\qquad N^{3/2}=(hq)^{15},\\
  B^{k-1}/3\le\tau\le B^{k-1}/2.
 \end{gathered}
\]
Consequently, if
\[
 \beta=1515k-\frac{k-1}{2}=\frac{3029k+1}{2},
\]
then
\[
 \sqrt{2}\,rB^\beta<a_r\le2^{15}\sqrt{3}\,rB^\beta.
\]
Since $a_r/2\le n_r\le a_r$ for sufficiently large $r$, there are fixed
constants $A_k,D_k>0$ such that
\begin{equation}\label{alt:cardinality-growth}
 A_kr^\alpha\le n_r\le D_kr^\alpha,
 \qquad \alpha=1+30\beta=45435k+16,
\end{equation}
for every sufficiently large prime $r$. In particular, put
\begin{equation}\label{alt:exponent}
 \eta=\frac2\alpha=\frac{2}{45435k+16}.
\end{equation}
The upper bound in \eqref{alt:cardinality-growth} gives
$r^2\ge D_k^{-\eta}n_r^\eta$. Combining this with
\eqref{alt:scaled-area}, we obtain
\begin{equation}\label{alt:prime-area}
 \Delta(P_r)\ge c_*n_r^{-2+\eta},
 \qquad c_*=\frac1{64D_k^\eta}>0.
\end{equation}
Thus the same cardinality estimate that will permit interpolation also
converts the factor $r^2$ into a fixed power of the number of points.

For a sufficiently large integer $n$, put
$m=\lceil(n/A_k)^{1/\alpha}\rceil$.
Lemma~\ref{norm:prime-interval} supplies a prime $r$ with $m<r\le2m$.
Once $(n/A_k)^{1/\alpha}\ge1$, we have
$m\le2(n/A_k)^{1/\alpha}$, so \eqref{alt:cardinality-growth} gives
\[
 n\le n_r\le C_kn,
 \qquad C_k=\max\left\{1,\frac{4^\alpha D_k}{A_k}\right\}.
\]
This uses only the two-sided cardinality bounds, not monotonicity of $n_r$.
Retain any $n$ points of $P_r$. Deleting points cannot decrease the minimum
triangle area. Since $0<\eta<2$, \eqref{alt:prime-area} therefore yields
\[
 \Delta(n)\ge c_* n_r^{-2+\eta}
 \ge c_*C_k^{-2+\eta}n^{-2+\eta}.
\]
Taking $c_1=c_*C_k^{-2+\eta}$ proves Theorem~\ref{intro:main}.
All constants are absolute because $k$ is fixed independently of $r$ and
$n$. The formula \eqref{alt:exponent} is explicit; no attempt has been
made to optimize the construction parameters.
\end{proof}
